% !TEX root = ../notes.tex

\documentclass[../notes.tex]{subfiles}

\begin{document}

\section{May 11}
Here we go.

\subsection{More on \texorpdfstring{$L$}{ L}-Functions}
Let's do an example.
\begin{exe}
	Fix a cusp form $f$ of $\op{SL}_2(\ZZ)$ and weight $k$. Then the completed $L$-function of $f$ is the zeta integral of the associated automorphic form.
\end{exe}
\begin{proof}
	Define the automorphic form $\varphi\colon\op{GL}_2(\QQ)\backslash\op{GL}_2(\AA_\QQ)\to\CC$ by
	\[\varphi\left(\begin{bmatrix}
		a & b \\ c & d
	\end{bmatrix}\right)\coloneqq(ci+d)^kf\left(\frac{ai+b}{ci+d}\right)\]
	for $\begin{bsmallmatrix}
		a & b \\ c & d
	\end{bsmallmatrix}\in\op{GL}_2(\RR)$; this fully defines $\varphi$ by strong approximation. Now, the zeta integral is
	\[Z(s,\varphi)=\int_{\QQ^\times\backslash\AA_\QQ^\times}\varphi\left(\op{diag}(a,1)\right)\left|a\right|^{s-1/2}\,d^\times a.\]
	Now, the integrand descends to a function on $\QQ^\times\backslash\AA_\QQ^\times/\widehat\ZZ^\times$, which is $\RR^+$, so
	\[Z(s,\varphi)=\int_{\RR^+}\varphi(\op{diag}(a,1))a^{s-1/2}\,\frac{da}a.\]
	Plugging into our definition of $\varphi$, we see that this is
	\[Z(s,\varphi)=\int_{\RR^+}f(ai)a^{s-1/2}\,\frac{da}a.\]
	Now, recall that $f$ admits a $q$-expansion $\sum_{n\ge1}a_nq^n$ where $q=e^{2\pi inz}$. Integrating term-by-term, we find that
	\[Z(s,\varphi)=(2\pi)^{-s}\Gamma(s)\sum_{n\ge1}\frac{a_n}{n^s},\]
	so we have recovered the completed $L$-function of the modular form $f$!
\end{proof}
\begin{remark}
	There is a functional equation
	\[\Lambda(s,f)=(-1)^{k/s}\Lambda(k-s,f).\]
	One can show this more directly using the invariance of $f$ under the Weyl element $\begin{bsmallmatrix}
		0 & 1 \\ -1 & 0
	\end{bsmallmatrix}$.
\end{remark}
A similar proof works in general.
\begin{proposition}
	Fix a global field $F$, and fix a cuspidal automorphic form $\varphi$ of $\op{GL}_2(\AA_F)$ of central character $\omega$. Then $w\coloneqq\begin{bsmallmatrix}
		0 & -1 \\ 1 & 0
	\end{bsmallmatrix}$ has
	\[Z(\varphi,\chi)=Z\left(R(w)\varphi,w^{-1}\chi^{-1}\right).\]
\end{proposition}
\begin{proof}
	Direct calculation.
\end{proof}
\begin{remark}
	Let $f$ be a cusp form of level $\op{SL}_2(\ZZ)$ and weight $k$, and let $\varphi$ be the associated automorphic form. Then the central character of $\varphi$ is $\left|\cdot\right|^{k}$. Thus, one can recover the functional equation of $\Lambda(s,f)$.
\end{remark}
We close our global discussion by trying to generalize.
\begin{definition}[Rankin--Selberg integral]
	Fix global field $F$ and automorphic representations $\pi$ and $\pi'$ of $\op{GL}_n$ and $\op{GL}_{n-1}$, respectively. Further assume that $\pi$ is cuspidal. Then for $\varphi\in\pi$ and $\varphi'\in\pi'$, we define
	\[Z(s,\varphi,\varphi')\coloneqq\int_{\op{GL}_{n-1}(F)\backslash\op{GL}_{n-1}(\AA_F)}\varphi\left(\begin{bmatrix}
		A \\ & 1
	\end{bmatrix}\right)\varphi'(A)\left|\det A\right|^s\,dA.\]
\end{definition}
\begin{remark}
	As before, one can check that this integral converges, admits an Euler product when $\varphi$ and $\varphi'$ are pure tensors, and there is a similar functional equation.
\end{remark}
\begin{remark}
	There is a further generalization to $\mathrm{GL}_n\times\mathrm{GL}_m$ for $n$ and $m$ arbitrary. It is slightly more complicated when $n=m$.
\end{remark}

\subsection{Local Zeta Integrals}
We now pass to the local story. For archimedean local fields, one just gets some products of $\Gamma$-functions, so we will ignore this and instead just treat the local case.
\begin{definition}[local integral]
	Fix a nonarchimedean local field $F$ and an irreducible automorphic representation $\pi$ of $\op{GL}_2(F)$. Further, suppose that $\pi$ admits a Whittaker functional $\Lambda_\pi$, and let $\mc W_\pi\subseteq C^\infty((N(F),\psi_N)\backslash\op{GL}_2(F))$ be the Whittaker model. Then for each $W\in\mc W_\pi$, we define the \textit{local zeta integral}
	\[Z(W,\chi)\coloneqq\int_{F^\times}W(\op{diag}(a,1))\chi(a)\,d^\times a.\]
\end{definition}
\begin{remark}
	As before, one can define
	\[Z(s,W,\chi)\coloneqq\int_{F^\times}W(\op{diag}(a,1))\chi(a)\left|a\right|^{s-1/2}\,d^\times a.\]
\end{remark}
\begin{remark}
	The function $s\mapsto Z(s,W,\chi)$ is a rational function in $\CC(q^{-s})$, where $q$ is the cardinality of the residue field of $F$. One can see this by a direct expansion: stratify the integral over $F^\times$ into $\bigsqcup_{r\in\ZZ}\varpi^r\OO_F^\times$, and then we can attempt to integrate over each piece.
\end{remark}
Here is how local integrals give rise to local $L$-factors.
\begin{proposition} \label{prop:get-local-l-factor}
	Fix a nonarchimedean local field $F$ and an irreducible automorphic representation $\pi$ of $\op{GL}_2(F)$. Suppose that $\pi$ admits a Whittaker model $\mc W_\pi$. Then the rational functions
	\[\{Z(s,w)\}_{W\in\mc W_\pi}\]
	admits a common denominator in $\CC[q^{-s}]$, and it can be normalized to have constant term $1$.
\end{proposition}
\begin{proof}
	Omitted.
\end{proof}
\begin{definition}
	Fix a nonarchimedean local field $F$ and an irreducible automorphic representation $\pi$ of $\op{GL}_2(F)$. Suppose that $\pi$ admits a Whittaker model $\mc W_\pi$. Then we define $L(s,\pi)$ to be the (reciprocal of the) common denominator in \Cref{prop:get-local-l-factor}.
\end{definition}
\begin{example}
	Suppose that $\pi$ is an irreducible spherical principal series representations $\op{Ind}_B^G(\chi_1,\chi_2)$. It turns out that
	\[L(s,\pi)=\frac1{\left(1-\alpha_1q^{-s}\right)\left(1-\alpha_2q^{-s}\right)}.\]
	Indeed, we computed the relevant integrals on the homework.
\end{example}
\begin{example}
	Similarly, one can compute that $L(s,\pi)=1/\left(1-\alpha q^{-s}\right)$ if $\pi$ is Steinberg, and $L(s,\pi)=1$ if $\pi$ is supercuspidal.
\end{example}
The last analytic property we want is a functional equation.
\begin{theorem}
	Fix a nonarchimedean local field $F$ and an irreducible automorphic representation $\pi$ of $\op{GL}_2(F)$ of central character $\omega$. Suppose that $\pi$ admits a Whittaker model $\mc W_\pi$. For $w\coloneqq\begin{bsmallmatrix}
		0 & -1 \\ 1 & 0
	\end{bsmallmatrix}$, the function
	\[\frac{Z\left(1-s,\pi(w)W,\omega^{-1}\chi^{-1}\right)}{Z(s,W,\chi)}\]
	is independent of $W$.
\end{theorem}
% \begin{proof}
% 	Omitted.
% \end{proof}
\begin{remark}
	This function may depend on $\pi$, $\chi$, $s$, and even on the character $\psi$ of $N(F)$ which gives rise to the Whittaker model.
\end{remark}

\subsection{The Kirillov Model}
The proofs of the various properties requires some preparation.
\begin{remark}
	A smooth representation $\pi$ of $\op{GL}_2(F)$ can be restricted to $F=N(F)$. The smoothness then induces an action by $C_c^\infty(F)$, which gains the structure of an algebra by convolution. By taking the Fourier transform, we see that $\pi$ attains an action by $C_c^\infty(F^*)$ where the multiplication is given by pointwise multiplication. (Here, $F^*$ is the Pontryagin dual!)
\end{remark}
Now, there is a ``universal'' such module over $C_c^\infty(F^*)$, namely $C_c^\infty(F^*)$. Here are some other such modules.
\begin{example} \label{ex:sheaf-to-cc-mod}
	For any sheaf $\mc E$ on the topological space $F^*$, the locally constant functions $C^\infty(F^*)$ act naturally by pointwise multiplication on $\Gamma(F^*,\mc E)$. By further restriction, we see that $C^\infty_c(F^*)$ acts naturally on $\Gamma_c(F^*,\mc E)$.
\end{example}
This discussion upgrades into an equivalence of categories.
\begin{proposition}
	Fix a totally disconnected locally compact space $X$. Then the functor $\Gamma_c(X;-)$ defines an equivalence
	\[\op{Sh}_\CC(X)\to\op{Mod}C_c^\infty(X).\]
\end{proposition}
\begin{proof}
	The forward functor was given in \Cref{ex:sheaf-to-cc-mod}. For the inverse functor, we start with a $C_c^\infty(X)$-module $M$, and we would like to define a sheaf $\mc E_M$ on $X$. Well, $X$ admits a basis given by compact open subsets $U$, so it is enough to merely define $\mc E_M(U)$, for which we use $\mc E_M(U)\coloneqq 1_U\cdot M$.
\end{proof}
Now, given a smooth representation $\pi$ of $F$, we may hope to understand the relevant sheaf $\mc E_\pi$. For example, one has the stalk
\[(\mc E_\pi)_0=\colim_n1_{\varpi^n\OO_F}\cdot V_\pi.\]
Now, the Fourier transform of $1_{\varpi^n\OO_F}$ is (up to scalar) $1_{\varpi^{-n}\OO_F}$ (possibly with some different contribution). Thus, we are interested in computing
\[(\mc E_\pi)_0=\colim_n1_{\varpi^{-n}\OO_F}*V_\pi,\]
so we see that we are basically taking a direct limit approaching the $N(F)$-invariants of $V_\pi$, which gives the co-invariants $(V_\pi)_{N(F)}$.
\begin{definition}[Jacquet module]
	Fix a nonarchimedean local field $F$ and a smooth representation $\pi$ of $\op{GL}_2(F)$. Then its \textit{Jacquet module} $J(\pi)$ is
	\[(V_\pi)_{N(F)}\coloneqq\frac{V_\pi}{\{\pi(u)-v:u\in N(F)\}}.\]
\end{definition}
\begin{remark}
	The preceding discussion shows that the associated sheaf $\mc E_\pi$ on $F$ has $(\mc E_\pi)_0$ equal to $(V_\pi)_{N(F)}$.
\end{remark}
\begin{remark}
	More generally, for $\psi\in F^*$, one finds that $(\mc E_\pi)_\psi=(V_\pi)_{(N(F),\psi)}$. For example, the dual is
	\[\op{Hom}((\mc E_\pi)_\psi,\CC)=\op{Hom}_{N(F)}(V_\pi,\CC_\psi),\]
	which is the space of Whittaker functionals for $\psi$. It follows that this stalk has dimension at most one.
\end{remark}
We thus have two cases.
\begin{itemize}
	\item If $\pi$ admits no Whittaker functional, then $\mc E_\pi$ is supported at $0\in F^*$. When $\pi$ is irreducible, this is only possible with $\pi$ is one-dimensional. Indeed, in this case, $\pi$ is isomorphic to its Jacquet module. We will argue shortly that $\dim J(\pi)\le2$, and $\dim J(\pi)=2$ only occurs for the irreducible principal series. Alternatively once we know that $V_\pi$ is finite-dimensional, one can find an $N(F)$-invariant vector $v$. By smoothness, $v$ is also fixed by a congruence subgroup, and this congruence subgroup with $N(F)$ generates all of $\op{SL}_2(F)$, so $\pi$ factors through the determinant.
	\item If $\pi$ admits a Whittaker functional, then the stalks of $\mc E_\pi$ are all one-dimensional away from zero. In fact, we claim that $\mc E_\pi|_{F^*\setminus\{0\}}\cong\CC$: indeed, $\mc E_\pi$ admits an $F^\times$-equivariant structure, so one can trivialize all stalks at once!
\end{itemize}
Let's now focus on the second case. The moral is that $V_\pi=\Gamma_c(F^*;\mc E_\pi)$ can now be computed as
\[0\to\Gamma_c(F^*\setminus\{0\};\mc E_\pi)\to\Gamma_c(F^*;\mc E_\pi)\to(\mc E_\pi)_0\to0,\]
where the last map is surjective because it is simply the surjection of $V_\pi$ onto its co-invariants. Notably, in the first case, $\Gamma_c(F^*\setminus\{0\};\mc E_\pi)=0$; in the second case, it is $C_c^\infty(F^*\setminus\{0\})$ with the natural $F^\times$-action.

Thus, the most mysterious remaining object is the Jacquet module.
\begin{remark}
	By Frobenius reciprocity, given a character $\chi$ of $T(F)$, we find that
	\[\op{Hom}_{T(F)}((V_\pi)_{N(F)},\CC_\chi)=\op{Hom}_{G(F)}\left(V_\pi,\op{Ind}_B^G\chi\right).\]
	Now, writing $\chi=(\chi_1,\chi_2)$, we see that such an embedding exists for at most two characters $\chi$, and the two characters are related by swapping them. It follows that $\dim J(\pi)\le2$.
\end{remark}
We can now separate our study of representations by $J(\pi)$.
\begin{example}
	One has that $\dim J(\pi)=2$ if and only if $\pi$ is an irreducible principal series $\op{Ind}_B^G\chi$. Explicitly, the induced map $V_\pi\to\op{Ind}_B^G\chi$ embeds $V_\pi$ into the principal series, and then one checks that the target is actually irreducible! Here, for the principal series to be irreducible, the character $\chi=(\chi_1,\chi_2)$ must have $\chi_1\chi_2^{-1}\notin\left\{\left|\cdot\right|^{\pm1}\right\}$.
\end{example}
\begin{example}
	One has that $\dim J(\pi)=1$ if and only if $\pi$ is one-dimensional or the Steinberg representation. Namely, this occurs when $\pi$ is found inside $\op{Ind}_B^G\chi$, where $\chi=(\chi_1,\chi_2)$ has $\chi_1\chi_2^{-1}\in\left\{\left|\cdot\right|^{\pm1}\right\}$.
\end{example}
\begin{definition}[supercuspidal]
	Fix a nonarchimedean local field $F$. Then $\pi$ is \textit{supercuspidal} if and only if $J(\pi)=0$.
\end{definition}
\begin{remark}
	Our above discussion shows that even the mysterious supercuspidal representations $\pi$ can be realized as the space $C_c^\infty(F^*\setminus\{0\})$.
\end{remark}
\begin{remark}[Kirillov embedding]
	In fact, one can typically embed into $C^\infty(F^*\setminus\{0\})$. Namely, if we have a Whittaker model $\mc W_{\pi,\psi}$, then one can send $W\in\mc W_{\pi,\psi}$ to the function
	\[\psi_a\mapsto W(\op{diag}(a,1))\]
	on $F^*\setminus\{0\}$. This induces a nonzero map from $V_\pi$ to $C^\infty(F^*\setminus\{0\})$. The moral is that $V_\pi$ is not too far away from $C_c^\infty(F^*\setminus\{0\})$, and on this space, $Z(W,s)$ actually converges. \Cref{prop:get-local-l-factor} follows if we actually keep track of what the Jacquet module can do to denominators.
\end{remark}
\begin{remark}
	Let's close our course with some further outlook.
	\begin{itemize}
		\item The next topic would be Eisenstein series. This grants a spectral decomposition of the space of automorphic forms.
		\item One can try to classify the supercuspidal representations explicitly.
		\item The construction of $L$-functions and their integral representations is an area of current research. Our best (conjectural) explanation for what is going on comes from the relative Langlands program (in the global case) and the Braverman--Kazhdan program (in the local case).
		\item This entire course has only lived on one side of the Langlands program. Studying the so-called motivic side and the relationships between them is also a current area of research.
	\end{itemize}
\end{remark}

\end{document}